\chapter[Tragedy, Ruins, and Storms]{Why Tragedy, Ruins, and Storms Can Fascinate Us}
\section{A Postcard from the Ruins}
First, pick up an object: a postcard.

Tourist gift shops sell many. Its image contains neither flowers nor smiling faces nor a new palace, but ruins: perhaps the Great Fountain remains at Beijing's Old Summer Palace. Broken white stone columns stand weatherworn, weeds pushing through cracks beneath an excessively clear blue sky. This is a wound left by fire and plunder, something broken.

Yet people buy it, display it on desks, and send it to friends. Some spend half an hour photographing it, declaring it wonderfully evocative.

Pause over the strangeness. Sellers could print reconstructions of an intact palace, gilded and exquisitely carved, without a flaw. They do, but those often sell less than broken columns. An intact building is "pretty"; a ruined one "overwhelming." Those differ, and the second seems more valuable.

Other evidence surrounds you. Cinema cities collapse and floods swallow streets while viewers happily eat popcorn, award high ratings, and recommend the film. Bestselling novels often kill someone or break hearts. Time-lapse videos of storm clouds advancing like walls draw millions of views. Horror-game streamers scream while chat fills with laughter, shouting "I'm done" while clicking "Continue."

We spend money and time actively seeking sadness, fear, and destruction. Why do tragedies, ruins, and storms, things we seemingly should avoid, fascinate us?

\section{Crying on a Hillside}
Enter a scene.

A spring morning in fifth-century BCE Athens, at an open-air theater below the Acropolis's southern slope. The Dionysian dramatic competition begins today. Before full daylight, thousands fill the stone tiers: older men with sticks, fathers holding children, craftsmen clutching coins, merchants visiting from abroad. Olive oil and dried figs scent the air; wine mixed in waterskins passes around. The sun rises from the sea, dazzling against whitewashed stone.

A signal sounds. The chorus enters from the side, masked actors moving to the center. A tragedy begins: a noble, fearless hero advances through an irreversible error toward destruction. You know the plot; audiences grew up with these stories, and the ending holds no surprise. Yet at the final lament, as fate crushes the hero like a millstone, the hillside becomes quiet enough to hear wind. The muscular man who joked loudly beside you on the way sits motionless, tears unheeded. Suppressed sobs begin in front and spread.

People leave red-eyed. Strangely, nobody complains or wants a refund. They discuss how this year's play cut deeper than last year's and how an actor's lines sent shivers through them. Days later, Athens rewards the playwright who made it suffer most. Next spring, people return, cry again, and reward the next writer who moves them most painfully.

Remember the hillside: thousands of sane, busy adults rise early, crowd onto stone, and sit in sun all day to watch fate destroy someone, weep, then judge the money well spent. This is no private eccentricity but a citywide ritual, part of an official festival and civic life. People debate the best play as we debate film awards. Athens openly placed sorrow at public life's center and rewarded it.

Nor was this an Athenian oddity. Two millennia later, you leave a cinema red-eyed and queue for tissues. You are the same person.

\section{Paying for Discomfort: A Genuine Paradox}
State the conflict before answering.

Ordinarily we call sadness and fear bad, pain something to avoid. Advice begins there: do not be sad, do not fear, look on the bright side. Someone charging fifty yuan for two hours of sadness with a helping of fear would sound mad.

Yet tragedy, horror, ruins, and disaster films sell exactly that service, successfully.

Philosophers call the puzzle the \textbf{paradox of tragedy}, or negative emotion. Three individually plausible statements cannot all hold:

\begin{itemize}
\item People always avoid things that make them sad or afraid.
\item Tragedy and horror genuinely produce sadness and fear; if they leave us untouched, we call them failures.
\item People actively seek tragedy and horror, sometimes liking them more the worse they feel.
\end{itemize}
At least one statement must be false. Which?

Perhaps the second: fictional fear is not real because we know it is fiction. Recall your racing heart and shielding hands. Perhaps the third: we merely tolerate sadness for plot and effects. Yet tearjerkers advertise making you cry as the attraction itself. Perhaps the first: people do not always avoid pain. Then a foundation of everyday advice, pain is bad, begins to loosen.

The paradox inhabits more than theaters: broken Summer Palace columns, storm time-lapses, frightening images you peek at through fingers. How do sadness, fear, destruction, and overwhelming force become enjoyment? Is humanity unwell, or have we misused "pleasure"?

Choose a provisional position: are hillside Athenians and cinema audiences healthy ordinary people, or collectively doing something somewhat unhealthy?

\section[Plato's Ban and His Student's Defense]{The Philosopher Who Banished Poets and His Student's Defense}
People have argued for two millennia. Give both positions their full strength.

\textbf{First: indulging sorrow and fear harms us and needs restraint.}

Plato is its strongest famous representative. In the Republic he broadly proposes excluding tragic poets from the ideal city. He was no ignorant outsider: he wrote beautifully and reportedly loved drama in youth. Knowing poetry's power made him fear it. Reducing his reasons to a stuffy spoilsport is unfair.

He had at least three concerns. Emotion becomes habitual and spills over. Weeping at others' disasters trains being led by feeling; indulgent theatergoers may collapse rather than act in real trouble. Second, imitation shapes people. Actors imitate madness and grief, spectators imitate emotions, and repeated imitation changes character. Immersing youths in uncontrolled sorrow teaches from bad textbooks. Third, theatrical emotion bypasses reason. Judgment needs calm; tragedy specializes in striking before calm is possible. A craft conquering audiences before they can think should not educate the city.

Modern versions are immediately recognizable: restricting children's exposure to horror and gore, film ratings, concern over disaster clips harvesting emotion. Plato's fear lives in content-classification systems.

\textbf{Second: sadness and fear are not harmful but necessary.}

Plato's student Aristotle defended tragedy not by denying pain but by asking what it \textbf{does}. Stories arouse pity and fear, releasing and exercising them. As bodies sweat, minds need outlets and rehearsals for accumulated feelings. Safe emotional use leaves viewers clearer and healthier, not crushed. We will examine his famous formulation later. For now, note the position: painful feelings are not the disease; suppression and nowhere to put them are.

\textbf{Another voice comes from China.} Confucius advocated neither banishing poetry nor unrestrained grief. The Analects, "Eight Rows," judges the Book of Songs' opening poem:

\begin{quote}
The Master said: "Guan Ju is joyful without excess, sorrowful without harm."
\end{quote}
Its joy does not overflow, its sorrow does not injure. This stands neither with Plato, since sorrow need not be expelled, nor with "the more pain the better," since excess harms. It adds a dimension: not whether, but how much.

Three voices: restrain, use freely, observe measure. You may instinctively favor one. Wait; we still need tools.

\section[Mapping What Moves Us]{A Bridge for Thought: Mapping What Moves Us}
Separate the overworked word "beauty." Landscapes, paintings, and ruins' "broken beauty" get one label; insufficient vocabulary confuses accounts. Aestheticians, studying why things move us, distinguish several experiences. Their map is portable.

For anything moving you, ask two questions.

\textbf{First: does it accommodate or overwhelm me?} Flowers, gentle streams, and lullabies comfort and invite closeness. This is \textbf{the beautiful}: small, soft, whole, unthreatening. Storms, stars, abysses, and tsunamis are immense, rough, overpowering. They do not please you; they diminish you. This is \textbf{the sublime}, sometimes called magnificent beauty.

\textbf{Second: does safety separate us?} Fear behind a screen differs from fear behind your own front door. A protective barrier may convert fear into enjoyment; without it, fear is simply disaster.

These questions yield a simple map:

\noindent
\begin{tabularx}{\linewidth}{lllX}
\toprule
Type & Typical object & Bodily response & Key condition \\
\midrule
Beautiful & \shortstack[l]{Streamside\\ flowers, lullabies} & \shortstack[l]{Relaxation,\\ attraction} & Small, harmless, harmonious \\
Sublime & \shortstack[l]{Storms, stars,\\ abysses} & \shortstack[l]{Overwhelmed, then\\ uplifted} & Vast and powerful, but cannot harm me now \\
Grotesque & \shortstack[l]{Melting clocks,\\ human-headed fish} & \shortstack[l]{Unease mixed with\\ laughter} & Familiar things wrongly combined \\
Horrific & \shortstack[l]{Dark footsteps,\\ screen monsters} & \shortstack[l]{Tension, urge to\\ flee} & Threat feels close, yet a barrier remains \\
\bottomrule
\end{tabularx}

\smallskip
The \textbf{grotesque} deserves more explanation. Its recipe is not size but wrongness: walking statues, human-faced fish, sinister songs in innocent children's voices. It makes us uneasy and amused by triggering recognition and mismatch together. It borders humor and horror: the same misassembled face becomes comedy in brighter light, nightmare in darkness.

Distinguish \textbf{horror} from the sublime too. Sublime objects overwhelm at a grand distance; storms do not target you individually. Horror watches, approaches, hides behind the door. Sublimity raises your head; horror makes you shrink. They are easily confused because one product often packages both: a disaster film's towering wave brings sublimity, followed by horror as something moves in a dark corridor.

Three instructions accompany the map. Boundaries move: a volcano is sublime to a safe observer and disastrous to residents beneath it. Labels map relationships, not objects alone. Second, classification is not ranking: sublime is not higher than beautiful, horror not lower than tragedy. Third, the map clarifies questions, not answers. Disputes about unhealthy fascination often conflate beauty and horror under one word.

Cultures have their own maps. Japanese mono no aware names being moved by passing things: cherry blossoms tighten the heart precisely because they will soon fall. Beauty and sorrow become one feeling. Ancient Indian dramatic theory in the Natyashastra distinguishes artistic "flavors," including karuna, the peculiar enjoyment of sorrow. Roughly two thousand years ago, theory already asked why sadness tastes good. With "sorrow without harm," these show a shared discovery: sadness can be tasted and seasoned; cultures debate the recipe.

\section{A Passage from Burke}
Read primary evidence. Eighteenth-century Irish thinker Edmund Burke intensely examined sublimity in A Philosophical Enquiry into the Origin of Our Ideas of the Sublime and Beautiful, published in 1757. Part I, section VII presents his central claim, rendered here from the Chinese paraphrase:

\begin{quote}
Whatever can awaken ideas of pain or danger---whatever is terrifying, concerns terrifying objects, or operates like terror---is a source of the sublime: it produces the strongest emotion the mind can experience.
\end{quote}
Read slowly. Burke offers not the commonplace that lovely things please, but \textbf{ideas of pain and danger themselves supply sublimity's raw material}. Terror, not gentleness, excites the mind most strongly.

He distinguishes actual pain and danger from their ideas. A burning home gives fear, not enjoyment. Watching a distant storm from safety or a hero's downfall in theater presents danger's \textbf{idea}, not danger itself, creating pleasurable shock rather than pure pain. Terror torments at close range but fascinates at an appropriate distance: fire on your body is catastrophe, safely observed fire magnificent. Burke even proposed physiology: terror tensions nerves, moderate tension exercises them, and absent actual injury turns strain into exhilaration. Tragedy and storms attract by simulating threat; near-danger feeds us.

The theory has limits, explored below. Keep two tools: terror as sublimity's material, and separate accounts for danger and its idea. Burke cornered this problem over thirty years before Kant, who entered through the door he opened.

\section{Why Tragedy Feels Good: Three Accounts}
Compare explanations of hillside tears and cinema heartbeats. Separate what happens, voluntary viewing and reported enjoyment; why, the contested explanation; participants' accounts, shaken Athenians or your "devastating" film; and judgments of whether it is good. We compare why.

\textbf{First: catharsis, emotions need draining.}

An Aristotelian account treats pity and fear as accumulating within, with tragedy a safe outlet. Crying can relieve; "let it out" is common experience. Tension before horror and helpless laughter afterward also resemble emptying. Yet if viewers drained themselves, repeated viewing should gradually become unnecessary. Enthusiasts seek a hundred-and-first tragedy after a hundred. Many horror viewers finish too excited to sleep, not relaxed. The pump seems unable to empty, perhaps producing more as it works.

\textbf{Second: distance transforms pain through safety.}

Burke and Hume offer versions. Hume's "Of Tragedy" broadly argues that real pain meets real pleasure from eloquence, structure, and convincing performance; artistic pleasure converts pain, like bitterness yielding sweetness in fermentation. Distance adds the pleasure of knowing danger is elsewhere and oneself safe. One film enjoyed at home may torment in a supposedly haunted house: the formula remains, distance changes. But \textbf{Roman arenas} challenge it. No screen, no merely imagined danger: real blood and deaths brought tens of thousands cheering. Knowledge of fiction cannot explain everything. Human excitement at real suffering must enter an honest account. Nor does distance explain those who cannot bear horror even on a screen: protective distance differs between people.

\textbf{Third: cognition and meaning, safely rehearsing life.}

Another Aristotelian clue is the desire to know through imitation and story. Tragedy and disaster become low-cost simulations of loss, betrayal, death, and extremity, practicing "what would I do?" and revealing what matters to us through whom we mourn. Horror exercises heartbeat, flight, and survival without real costs. This explains thoughtful silence after great tragedy, not emptying but reflection, and lingering at ruins. Broken columns silently teach that even splendor falls, a lesson ignored in words but heard in stone. Yet this makes audiences too philosophical. Crude horror and disaster films sell without deep meaning, and why rehearse identically dozens of times? Who schedules fire drills three times weekly?

Each explains part: catharsis relief, distance safety's role, meaning lingering aftereffects. The arena stands beside them, reminding us that more primitive, less honorable pleasures may mingle. Acknowledging this is honesty, not defeat.

\section{A Deeper Challenge: Reason Takes over Fear}
Now the weightiest theory, Kant's sublime in the Critique of Judgment of 1790. His interest in destructive nature had historical roots. Lisbon's earthquake in 1755 destroyed the city and killed tens of thousands, making Europe question senseless natural power. Young Kant wrote essays seeking natural causes rather than divine punishment. Sublimity arose not as a study-room game but in an earthquake-frightened century.

Kant distinguished two ways of being overwhelmed.

\textbf{First: the mathematical sublime, beyond imagination's scale.} At the Grand Canyon or beneath the Milky Way, you try to contain the full dimensions mentally and fail. Imagination has limits; the object seems limitless. Yet when imagination concedes, reason arrives. You cannot picture the scale, but can \textbf{think} infinity. Concepts reach beyond senses. Frustration becomes exhilaration: immense nature itself can be contemplated through reason's concept of infinity. Sublimity is the self-respect of discovering thought beyond every sensory object.

\textbf{Second: the dynamical sublime, powerful enough to destroy me, but not now.} Storms, volcanoes, raging seas truly could crush us. From safety, fear yields another awareness: nature can overwhelm our strength but not demand our submission. Bodies may perish; judgment, dignity, and morality lie outside its jurisdiction. Sublimity, Kant broadly argues, inhabits us, not the storm. Being pressed down and raising our heads reveals another kind of height within.

Weigh its power. It explains the strange \textbf{uplift} beyond Burke's relief at survival. For Burke, storm-watching exercises nerves; for Kant, reason reviews its strength. One explains heartbeats, the other why someone descends a mountain feeling taller.

Weigh its boundaries too. Infinity and rational dignity demand abstraction absent from most screaming horror audiences. Kant illuminates the sublime, not the whole map: grotesque laughter and compulsive horror have little place. Bright theoretical boundaries make the darkness outside clearer. No single theory should be expected to explain all human fascination.

\section{Ruins as Photo Stops: Today's Ethical Boundary}
The ancient question plays daily on your screen, upgraded.

Fiction accelerates. Disaster-film cities fall repeatedly each year; horror games engineer fright through calculated jump scares and tuned dark sounds. Streamer screams become laughing compilations. One-minute tragic videos cue music and synchronized comments of emotional collapse. Standardized emotion achieves efficiency beyond Plato's imagination.

Reality becomes more difficult. Consider scenes existing today:

\begin{itemize}
\item Ruin tourism. Chernobyl's exclusion zone became popular, with queues for selfies at an abandoned Ferris wheel and increased visits after the television series. This was a real nuclear disaster: people died, others permanently lost homes, and consequences continue.
\item Disaster clips. Earthquake, flood, and fire footage spreads with sensational titles and advertising revenue. The filmer stood somewhere safe holding a phone.
\item True crime. Podcasts and videos turn actual murders into bedtime stories, the most successful reaching millions. Families sometimes unexpectedly encounter a relative's death presented as entertainment.
\end{itemize}
Here lies a line Burke and Kant did not draw, because the eighteenth century lacked this problem: the ethical boundary between \textbf{appreciating disaster-themed art} and \textbf{consuming real suffering}. Fiction injures nobody, permitting enjoyment long accepted across cultures. Real suffering connects fascination to someone else's fate. Lu Xun's "Medicine" offers a famous image: people buy buns soaked in an executed person's blood, believing them curative. Buyers need not be malicious; they numbingly consume death as something normal. The image still asks not whether you are kind but what your pleasure stands on.

Where is the boundary? Rather than ready answers, use three questions.

First, \textbf{fiction or documentary?} Are sufferers actors or real people? This determines whether viewing is aesthetic contemplation or spectatorship.

Second, \textbf{do those involved consent?} Ruins can be photographed, but lives once lived there belonged to people. Do former Chernobyl residents' feelings count when visitors pose jokingly before their childhood homes?

Third, \textbf{what does your viewing feed?} One disaster clip may bring needed attention or merely add a view and advertising income. Viewing has consequences, revenue shares, beneficiaries; it is not neutral.

These questions distinguish historical documentation from posed attention-seeking, understanding justice from seeking bloody bedtime detail. It is still our map, but beyond the safety separating viewer and object now stand real people.

\section[The Second You Raise Your Phone]{Your Question: The Second You Raise Your Phone}
Return to the postcard.

The Summer Palace's broken columns have stood for over a century. Photographers change; stones cannot object or feel offended, making photography seem uncomplicated. But suppose your phone faces not old ruins but a disaster happening now. You stand safely while real fire, panic, and people needing help fill the screen. Posting could inform many and gain followers. Putting it down might let you help, or at least prevent your pleasure from resting on somebody's worst day.

Raise it or not?

Weigh what you learned. Fascination with terror and destruction is ordinary, not illness. Documentation matters; history survives through records. Requiring every witness to rescue is too demanding. Conversely, arena cheers show liking does not establish rightness. Real suffering has an ethical boundary occupied by people unable to refuse. "Everybody does it" never proves justification.

At that moment, is the phone-holder witness, recorder, onlooker, or consumer? Are these genuine differences, or stories afterward used for self-persuasion?

There is no standard answer. But next time tragedy makes you cry, a storm overwhelms you, or a disaster clip fixes your gaze, you can see two things at once: your ancient machinery for fear and sorrow at work, and whether its object is fiction or another person's actual life.
